\chapter{Research: Existing File System Architecture}

\section{The Android Filesystem Foundation}

Android runs on the Linux kernel and uses a hierarchical tree structure for all files. Every file and directory is represented internally by an \textbf{inode} (index node) --- a 256-byte on-disk structure containing all metadata about a file except its name.

\subsection{The inode}

An inode stores: file size, owner UID and GID, permissions, three timestamps (created, accessed, modified), and a pointer to where the file's content blocks reside on disk. The file's name is stored separately in a directory entry that maps the name to the inode number. Renaming a file changes only the directory entry --- the inode is untouched, which is why renaming is instantaneous at the OS level.

\subsection{ext4: The Extent B-Tree}

Inside every inode, a 60-byte field called \texttt{i\_block} holds the root of an \textbf{Extent B-Tree} --- a balanced tree mapping logical file offsets to physical disk block addresses. For small files the tree fits entirely inside the inode (\textit{inline data}). For large files it expands to additional disk blocks. B-tree search, insert, and delete all execute in \bigO{\log n}$.

\subsection{Directory Indexing: HTree}

Small directories use a linear linked list of entries --- \bigO{n} for lookup. When a directory grows large, ext4 converts it to an \textbf{HTree} (Hashed B-Tree): filenames are hashed, and the hash becomes the B-tree key, enabling \bigO{\log n} directory lookup.

\begin{insightbox}[title={Critical Observation}]
Even at the kernel level, large directories use B-trees. Yet file explorer applications sitting above this foundation perform linear scans across thousands of files. The inefficiency is entirely in the application layer, not the kernel.
\end{insightbox}

\section{Android MediaStore}

File explorer apps cannot access ext4 directly --- Android's permission model prevents it. Instead, they query \textbf{MediaStore}: an SQLite database maintained by the OS indexing all media files.

MediaStore has four tables: \texttt{Images}, \texttt{Video}, \texttt{Audio}, and \texttt{Files}. SQLite uses B-trees for its column indexes. Indexed queries (sort by \texttt{date\_added}) execute in \bigO{\log n + k}. But substring filename search issues a \texttt{LIKE '\%query\%'} statement which cannot use any B-tree index, degrading to \bigO{n}.

\subsection{The Staleness Problem}

MediaStore is updated asynchronously by the Media Scanner. A file deleted by one application may continue to appear in MediaStore for seconds to hours. Samsung My Files further compounds this by maintaining its own SQLite database at \texttt{data/data/com.sec.android.app.myfiles/databases/myfiles.db} --- a separate cache that independently goes stale.

\section{Analysis of Major Competitors}

\begin{longtable}{@{}p{30mm}p{35mm}p{35mm}p{40mm}@{}}
\toprule
\textbf{App} & \textbf{Primary Index} & \textbf{Search Method} & \textbf{Critical Weakness} \\
\midrule
Google Files & MediaStore SQLite & SQL \texttt{LIKE} full scan & No substring index; stale data \\
Samsung My Files & MediaStore + myfiles.db & SQL \texttt{LIKE} + dir scan & Dual stale caches; permission gaps \\
Google Photos & Cloud vector DB + SQLite & Semantic vector similarity & Cloud-dependent; photos only \\
Windows Explorer & NTFS MFT + inverted index & Inverted word index & Slow first-run indexing \\
\bottomrule
\end{longtable}

Google Photos is the instructive outlier: it uses ML-generated vector embeddings for every photo and performs nearest-neighbour lookup in embedding space, enabling search for ``beach sunset'' without those words in any filename. The limitation is that this infrastructure runs in Google's cloud. Bringing equivalent capability fully on-device, for all file types, is FLUX's central innovation.

\chapter{The FLUX Index Architecture}

\section{Design Manifesto}

\begin{lawbox}
THE FLUX LAW: Every operation a user can trigger --- lookup, delete, search, insert, filter, list --- must execute in O(1) or O(log n) time. No full scans. No exceptions. If a data structure cannot guarantee this, we use a different one.
\end{lawbox}

\section{The File ID System}

Every file is assigned a unique 64-bit integer \textbf{FID} (File Identifier) at first discovery. All nine indexes reference the FID, never the path string.

\textbf{Why FIDs matter:}
\begin{itemize}
  \item \textbf{Rename safety:} Renaming or moving a file updates only the path index. All other nine indexes hold FIDs --- zero cascading updates.
  \item \textbf{Memory efficiency:} An FID is 8 bytes. A typical file path is 60--120 bytes. Posting lists using FIDs use 87--93\% less memory.
  \item \textbf{Bitset operations:} Integer FIDs enable set intersections via 64-bit CPU word operations --- AND, OR, NOT on millions of entries in microseconds.
  \item \textbf{\bigO{1} record access:} The master record for any file is one array index operation.
\end{itemize}

FIDs are allocated by \texttt{AtomicLong.incrementAndGet()} --- thread-safe and monotonically increasing. FID 0 is reserved as a null sentinel. Deleted FIDs enter a free pool only during scheduled background compaction.

\section{The Master Record Array}

\begin{lstlisting}[language=Kotlin, caption=Master Record Array]
val masterIndex = Array<FileRecord>(MAX_FILES) { FileRecord.EMPTY }

// O(1) lookup -- always, regardless of total file count
val record: FileRecord = masterIndex[fid]
\end{lstlisting}

The \texttt{FileRecord} is a 64-byte struct optimised for mobile RAM constraints:

\begin{longtable}{@{}p{30mm}p{20mm}p{12mm}p{68mm}@{}}
\toprule
\textbf{Field} & \textbf{Type} & \textbf{Size} & \textbf{Purpose} \\
\midrule
\texttt{fid} & uint64 & 8\,B & Unique file identifier \\
\texttt{parentDirFid} & uint64 & 8\,B & FID of parent directory \\
\texttt{nameOffset} & uint24 & 3\,B & Offset into shared name string pool \\
\texttt{nameLen} & uint16 & 2\,B & Byte length of filename \\
\texttt{pathOffset} & uint24 & 3\,B & Offset into shared path string pool \\
\texttt{pathLen} & uint16 & 2\,B & Byte length of full path \\
\texttt{size} & uint64 & 8\,B & File size in bytes \\
\texttt{mtime} & uint32 & 4\,B & Modified time (seconds since 2020-01-01) \\
\texttt{atime} & uint32 & 4\,B & Accessed time \\
\texttt{ctime} & uint32 & 4\,B & Created time \\
\texttt{mimeType} & uint16 & 2\,B & Index into global MIME type table \\
\texttt{flags} & uint32 & 4\,B & Deleted, hidden, pinned, indexed, starred \\
\texttt{vectorSlot} & uint24 & 3\,B & Slot in HNSW vector store \\
\texttt{accessCount} & uint16 & 2\,B & Open count (capped at 65,535) \\
\texttt{checksum} & uint64 & 8\,B & xxHash64 of content (deduplication) \\
\midrule
\textbf{Total} & & \textbf{64\,B} & \\
\bottomrule
\end{longtable}

\textbf{RAM at scale:} 100,000 files (typical) = 6.4\,MB. 500,000 files (power user) = 32\,MB. Fits entirely within the 50--80\,MB HOT cache tier.

\section{The Nine Indexes}

\begin{longtable}{@{}p{6mm}p{30mm}p{34mm}p{34mm}p{30mm}@{}}
\toprule
\textbf{\#} & \textbf{Name} & \textbf{Structure} & \textbf{Primary Query} & \textbf{Complexity} \\
\midrule
1 & Path Map & HashMap (xxHash64) & Exact path lookup & \bigO{1} \\
2 & Name Trie & Radix Trie & Prefix / autocomplete & \bigO{k} \\
3 & Token Index & HashMap + RoaringBitmap & Keyword search, AND/OR & \bigO{1}/\bigO{N/64} \\
4 & Directory Index & HashMap + sorted int[] & Folder listing & \bigO{1} \\
5 & Type Buckets & HashMap + RoaringBitmap & Filter by MIME type & \bigO{1} \\
6 & Size Index & Van Emde Boas Tree & Range by file size & \bigO{\log\log U} \\
7 & Time Index & Van Emde Boas Tree & Range by date & \bigO{\log\log U} \\
8 & Checksum Map & HashMap (xxHash64) & Duplicate detection & \bigO{1} \\
9 & HNSW Vector Graph & Multi-layer proximity graph & Semantic AI search & \bigO{\log n} \\
\bottomrule
\end{longtable}

\subsection{The Deletion BitSet --- Making Deletion O(1)}

\begin{lstlisting}[language=Kotlin, caption=O(1) Logical Deletion]
// User deletes N files -- O(1) per file for user-perceived latency
fun logicalDelete(fid: Int) {
    deletionSet.add(fid)                   // flip one bit
    masterIndex[fid].flags = masterIndex[fid].flags or FLAG_DELETED
    wal.append(WAL_DELETE, fid)            // crash-safe write-ahead log
}

// Every query result is filtered through:
val result = queryResult.andNot(deletionSet)   // O(N/64)

// Physical cleanup runs asynchronously during idle
fun backgroundCompaction() {
    deletionSet.forEach { fid -> physicallyRemoveFromAllIndexes(fid) }
    deletionSet.clear()
}
\end{lstlisting}

Deleting 10,000 files: user-perceived latency = \bigO{10{,}000/64} = 156 bitmap operations, under 2\,ms. The phone never freezes.

\section{Index 1 --- Path HashMap: O(1) Exact Lookup}

\begin{lstlisting}[language=Kotlin, caption=Path HashMap]
val pathMap = HashMap<Long, Int>()   // xxHash64(path) -> FID

fun lookup(path: String): FileRecord? {
    val hash = xxHash64(path.lowercase())       // ~5 nanoseconds
    val fid  = pathMap[hash] ?: return null    // ~10 nanoseconds
    return if (deletionSet.contains(fid)) null else masterIndex[fid]  // ~10 ns
}
// Total: ~25 nanoseconds. Identical for 10 files or 10,000,000.
\end{lstlisting}

xxHash64 runs at approximately 10\,GB/s on modern hardware. The full lookup sequence completes in $\sim$25 nanoseconds, independent of total file count.

\section{Index 2 --- Radix Trie: O(k) Prefix Search}

The Radix Trie (compressed prefix tree) enables autocomplete in \bigO{k} time where \textit{k} is the query prefix length --- regardless of how many million files are indexed.

FLUX maintains two separate tries:
\begin{itemize}
  \item \textbf{\texttt{nameTrie}} --- indexed on complete filenames (autocomplete)
  \item \textbf{\texttt{tokenTrie}} --- indexed on extracted tokens, enabling mid-word search
\end{itemize}

\textbf{Tokenisation:} \texttt{Q3\_Budget\_2025\_FINAL\_v2.xlsx} becomes \texttt{[q3, budget, 2025, final, v2, xlsx]} by splitting on underscores, hyphens, dots, spaces, CamelCase boundaries, and number-letter transitions.

\section{Index 3 --- Token Index + RoaringBitmap}

\begin{lstlisting}[language=Kotlin, caption=Multi-keyword AND search]
val tokenIndex = HashMap<String, RoaringBitmap>()

// Multi-keyword AND: O(N/64) via bitmap intersection
fun searchAll(vararg tokens: String): RoaringBitmap {
    var result = tokenIndex[tokens[0]] ?: return RoaringBitmap()
    for (i in 1 until tokens.size) {
        val bitmap = tokenIndex[tokens[i]] ?: return RoaringBitmap()
        result = RoaringBitmap.and(result, bitmap)
    }
    return result.andNot(deletionSet)
}
\end{lstlisting}

\textbf{Why RoaringBitmap over a plain bitset:}

\begin{longtable}{@{}p{40mm}p{30mm}p{38mm}p{28mm}@{}}
\toprule
\textbf{Property} & \textbf{Plain Bitset} & \textbf{RoaringBitmap} & \textbf{Winner} \\
\midrule
Memory (1\,M files, 1\% density) & 125\,KB fixed & $\approx$5\,KB & Roaring 25$\times$ \\
AND speed & \bigO{N/64} & \bigO{N/64} + skips zeros & Roaring 2--5$\times$ \\
Sparse tokens (rare words) & 125\,KB wasted & ArrayContainer: tiny & Roaring \\
Serialisation & Large & Compressed & Roaring \\
\bottomrule
\end{longtable}

Content indexing works identically: for text files (PDF, DOCX, TXT), the first 10,000 characters are tokenised and stored under a \texttt{content:} namespace in the same \texttt{tokenIndex}. Searching inside documents is the same \bigO{1} lookup as searching filenames.

\section{Indexes 4--8 --- Directory, Type, Size, Time, Checksum}

\textbf{Index 4 --- Directory Index:} \texttt{HashMap<Int, IntArray>} keyed on directory FID. Listing a folder: \bigO{1}.

\textbf{Index 5 --- Type Buckets:} \texttt{HashMap<String, RoaringBitmap>} by MIME type. Filtering all images: \texttt{typeBuckets["image/*"]} --- \bigO{1}.

\textbf{Index 6 --- Size Index (Van Emde Boas Tree):} Universe size $U = 2^{40}$ (1\,TB max). Range queries (``files between 10\,MB and 100\,MB'') in \bigO{\log\log U} --- approximately 13 operations regardless of file count.

\textbf{Index 7 --- Time Index:} Same Van Emde Boas Tree on timestamps. ``Files from last week'' in \bigO{\log\log U}$.

\textbf{Index 8 --- Checksum Map:} \texttt{HashMap<Long, IntArray>} keyed on \texttt{xxHash64(content)}. Duplicate detection in \bigO{1}. Computing xxHash64 of a 10\,MB file takes approximately 8\,ms on a modern Android SoC.

\section{Index 9 --- HNSW Vector Graph: O(log n) Semantic AI Search}

\subsection{Vector Embeddings}

Every file is passed through an on-device ML model (MiniLM-L6, 22\,MB, 384 dimensions) that converts the filename, path context, and document preview into a 384-dimensional numeric vector. Files with semantically similar content produce vectors that are mathematically close in this space.

\subsection{HNSW: Hierarchical Navigable Small World}

Brute-force nearest-neighbour search on 1,000,000 vectors requires 1,000,000 dot products. HNSW reduces this to approximately \bigO{\log n}:

\begin{itemize}
  \item \textbf{Top layer:} Few nodes, long-range connections --- coarse navigation
  \item \textbf{Middle layers:} More nodes, shorter connections
  \item \textbf{Bottom layer:} All nodes, short-range connections --- fine-grained recall
\end{itemize}

\begin{longtable}{@{}p{52mm}p{32mm}p{26mm}p{26mm}@{}}
\toprule
\textbf{Operation} & \textbf{Brute Force} & \textbf{HNSW} & \textbf{Speedup} \\
\midrule
Top-10 semantic search (1\,M files) & \bigO{n} = 1\,M ops & \bigO{\log n} $\approx$ 20 ops & $\sim$50,000$\times$ \\
Add file to graph & --- & \bigO{M \cdot \log n} & Sub-ms amortised \\
Delete from graph & --- & \bigO{M \cdot \log n} & Sub-ms amortised \\
\bottomrule
\end{longtable}

Recall accuracy: 95--98\% (approximate nearest neighbours).

\subsection{On-Device Model Selection}

\begin{longtable}{@{}p{44mm}p{22mm}p{18mm}p{18mm}p{38mm}@{}}
\toprule
\textbf{Model} & \textbf{Dimensions} & \textbf{Size} & \textbf{Latency} & \textbf{Use Case} \\
\midrule
MiniLM-L6 (recommended) & 384 & 22\,MB & $\sim$15\,ms & Default --- best balance \\
all-MiniLM-L3-v2 & 384 & 17\,MB & $\sim$8\,ms & Low-end devices \\
BERT-tiny & 128 & 5\,MB & $\sim$3\,ms & Ultra-budget devices \\
\bottomrule
\end{longtable}

Embedding generation runs exclusively when the device is charging, the screen is off, and the thermal state is COOL --- zero battery or thermal impact on normal usage.

\chapter{Mobile-First Architecture: Flutter + Native}

\section{The Split: Flutter vs Native Kotlin}

\begin{lawbox}
THE BRIDGE RULE: Native owns everything that touches hardware, the filesystem, or the OS scheduler. Flutter owns everything the user sees and touches. The bridge carries only FIDs and small metadata --- never raw bytes, never large strings.
\end{lawbox}

\begin{longtable}{@{}p{52mm}p{22mm}p{22mm}p{40mm}@{}}
\toprule
\textbf{Responsibility} & \textbf{Layer} & \textbf{Language} & \textbf{Rationale} \\
\midrule
FLUX Index (all 9 structures) & Native & Kotlin & Direct JVM memory, no Dart GC pressure \\
FileObserver / inotify & Native & Kotlin & OS API unavailable from Dart \\
Background indexing service & Native & Kotlin WorkManager & Survives app backgrounding \\
Thumbnail generation & Native & Kotlin ThumbnailUtils & Hardware-accelerated JPEG decoder \\
HNSW vector computation & Native & Kotlin + ONNX Runtime & NDK C++ inference \\
File read / write / delete & Native & Kotlin & Direct POSIX calls \\
All UI screens & Flutter & Dart & Single codebase, Android + iOS \\
State management & Flutter & Dart (Riverpod) & Reactive provider tree \\
\bottomrule
\end{longtable}

\section{RAM Tiering}

\begin{longtable}{@{}p{16mm}p{24mm}p{52mm}p{18mm}p{24mm}@{}}
\toprule
\textbf{Tier} & \textbf{Name} & \textbf{Contents} & \textbf{Max RAM} & \textbf{Eviction} \\
\midrule
TIER 1 & HOT CACHE & masterIndex, pathMap, dirIndex, tokenIndex, typeBuckets, deletionSet, ring buffer & 50--80\,MB & Never (foreground) \\
TIER 2 & WARM STORE & sizeVEB, timeVEB, checksumMap, nameTrie (full) & 80--150\,MB & On low memory \\
TIER 3 & COLD DISK & HNSW graph, content tokens, string pool & 600\,MB+ on disk & OS-managed mmap \\
\bottomrule
\end{longtable}

\begin{lstlisting}[language=Kotlin, caption=Memory Pressure Response]
override fun onTrimMemory(level: Int) {
    when (level) {
        TRIM_MEMORY_RUNNING_MODERATE -> {
            checksumMap.evictLRU(keepCount = 50_000)
        }
        TRIM_MEMORY_RUNNING_CRITICAL -> {
            warmStore.evictAll()
            tokenIndex.values.forEach { it.runOptimize() }
        }
        TRIM_MEMORY_BACKGROUND -> {
            wal.flush()
            indexingJob.pause()
        }
        TRIM_MEMORY_COMPLETE -> {
            fluxIndex.emergencyFlush()
        }
    }
}
\end{lstlisting}

\section{Thermal Management}

\begin{lstlisting}[language=Kotlin, caption=ThermalGovernor.kt]
class ThermalGovernor(private val pm: PowerManager) {
    enum class ThermalState { COOL, WARM, HOT, CRITICAL }

    fun currentState(): ThermalState = when {
        Build.VERSION.SDK_INT >= 29 -> when (pm.currentThermalStatus) {
            PowerManager.THERMAL_STATUS_NONE,
            PowerManager.THERMAL_STATUS_LIGHT    -> ThermalState.COOL
            PowerManager.THERMAL_STATUS_MODERATE -> ThermalState.WARM
            PowerManager.THERMAL_STATUS_SEVERE   -> ThermalState.HOT
            else                                 -> ThermalState.CRITICAL
        }
        else -> readCpuTempFallback()
    }

    fun workerParams(state: ThermalState) = when (state) {
        COOL     -> WorkerParams(threads=4, batchSize=500, delayMs=0)
        WARM     -> WorkerParams(threads=2, batchSize=200, delayMs=50)
        HOT      -> WorkerParams(threads=1, batchSize=50,  delayMs=200)
        CRITICAL -> WorkerParams(threads=0, batchSize=0,   delayMs=-1)
    }
}
\end{lstlisting}

\begin{longtable}{@{}p{48mm}p{24mm}p{18mm}p{16mm}p{30mm}@{}}
\toprule
\textbf{Condition} & \textbf{Indexing?} & \textbf{Threads} & \textbf{Batch} & \textbf{Rationale} \\
\midrule
Screen OFF + Charging & Full speed & 4 & 1,000 & Best window \\
Screen OFF + Battery & Moderate & 2 & 200 & Conservative \\
Screen ON + COOL & Background & 1 & 100 & No UI impact \\
Screen ON + WARM & Paused & 0 & 0 & Prevent heating \\
Screen ON + HOT & Paused & 0 & 0 & UX priority \\
Battery $<$ 20\% & Paused & 0 & 0 & Battery preservation \\
HNSW embedding & Charge only & 1 NPU & 200/10\,min & Most expensive task \\
\bottomrule
\end{longtable}

\chapter{Cross-Application Deletion Sync}

When WhatsApp deletes a sent photo, when Google Photos removes a duplicate, when Samsung My Files reorganises --- FLUX must know immediately. A file appearing in search results that does not exist on disk is a ghost entry. Ghost entries destroy user trust.

\section{4-Layer Defence System}

\begin{longtable}{@{}p{18mm}p{32mm}p{22mm}p{22mm}p{38mm}@{}}
\toprule
\textbf{Layer} & \textbf{Mechanism} & \textbf{Speed} & \textbf{Coverage} & \textbf{Complexity} \\
\midrule
Layer 1 & FileObserver (inotify) & $<$1\,ms kernel event & All watched directories & \bigO{1} per event \\
Layer 2 & MediaStore ContentObserver & 100\,ms--2\,s OS broadcast & All media files & \bigO{delta} \\
Layer 3 & Root observer on /sdcard & $<$1\,ms for new dirs & New unregistered dirs & \bigO{1} per new dir \\
Layer 4 & Delta reconciliation & 15-min idle cycle & Catch-all & \bigO{delta} not \bigO{n} \\
\bottomrule
\end{longtable}

\begin{insightbox}[title={The Kernel Guarantee}]
inotify is a Linux kernel subsystem. When any process --- including another app --- deletes a file, the kernel fires the inotify event within the same syscall that performed the deletion. FLUX is notified \textbf{before the deleting application's own UI updates}. Detection latency: $<$1 millisecond.
\end{insightbox}

\begin{lstlisting}[language=Kotlin, caption=Self-Registering FileObserverHub]
class FileObserverHub(private val flux: FluxIndex) {
    private val activeObservers = ConcurrentHashMap<String, FileObserver>()

    fun register(dirPath: String) {
        if (activeObservers.containsKey(dirPath)) return

        val observer = object : FileObserver(dirPath,
            CREATE or DELETE or MOVED_FROM or MOVED_TO or CLOSE_WRITE) {
            override fun onEvent(event: Int, filename: String?) {
                filename ?: return
                val fullPath = "$dirPath/$filename"
                when (event and ALL_EVENTS) {
                    CREATE -> {
                        val f = File(fullPath)
                        if (f.isDirectory) {
                            register(fullPath)               // self-register new dirs
                            flux.scanDirAsync(fullPath)
                        } else {
                            flux.insertAsync(fullPath)
                        }
                    }
                    DELETE, MOVED_FROM -> {
                        flux.logicalDelete(fullPath)         // O(1) tombstone
                        activeObservers.remove(fullPath)?.stopWatching()
                    }
                    MOVED_TO  -> flux.insertAsync(fullPath)
                    CLOSE_WRITE -> flux.invalidateChecksumAndThumb(fullPath)
                }
            }
        }
        observer.startWatching()
        activeObservers[dirPath] = observer
    }
}
\end{lstlisting}

\chapter{Data Schema and Persistence}

\section{Persistence Architecture}

FLUX uses a Write-Ahead Log (WAL) for crash safety and memory-mapped files for efficient cold-start recovery.

\begin{longtable}{@{}p{42mm}p{28mm}p{68mm}@{}}
\toprule
\textbf{File} & \textbf{Format} & \textbf{Contents} \\
\midrule
\texttt{flux.wal} & Binary sequential & Every insert/delete/update since last checkpoint \\
\texttt{master.bin} & Binary array & Serialised FileRecord[] --- mmap on startup \\
\texttt{pathmap.bin} & Binary & Serialised HashMap (xxHash64 to FID) \\
\texttt{tokenindex.bin} & Binary & Serialised HashMap + RoaringBitmaps \\
\texttt{dirindex.bin} & Binary & Serialised directory children arrays \\
\texttt{typebuckets.bin} & Binary & Serialised MIME type bitmaps \\
\texttt{deletion.bin} & Binary & Serialised deletionSet RoaringBitmap \\
\texttt{hnsw.bin} & Binary & HNSW graph layers (memory-mapped, 600\,MB+) \\
\texttt{strings.pool} & Binary & String pool byte array (mmap) \\
\texttt{trash.log} & Binary & Tombstoned FIDs + deletion timestamps \\
\texttt{settings.json} & JSON & User preferences \\
\bottomrule
\end{longtable}

\section{WAL Entry Format}

\begin{lstlisting}[language=Kotlin, caption=WAL Entry -- 32 bytes]
data class WalEntry(
    val magic:     UInt,   // 4B -- 0xFLUX for validation
    val sequence:  Long,   // 8B -- monotonic sequence number
    val timestamp: Long,   // 8B -- epoch milliseconds
    val opCode:    Byte,   // 1B -- INSERT=1, DELETE=2, UPDATE=3, RENAME=4
    val fid:       Int,    // 4B -- target file ID
    val payload:   Short,  // 2B -- operation-specific data
    val checksum:  Int,    // 4B -- CRC32 of entry for corruption detection
    // padding    1B       -- alignment to 32 bytes
)
\end{lstlisting}

\section{Settings Schema}

\begin{lstlisting}[caption=settings.json schema]
{
  "index": {
    "thermal_throttle": true,
    "embed_on_wifi_only": false,
    "max_watch_depth": 2,
    "content_index_enabled": true,
    "content_index_max_file_size_mb": 10
  },
  "ui": {
    "theme": "system",       // "light" | "dark" | "system"
    "default_view": "list",  // "list" | "grid"
    "grid_columns": 3,
    "show_hidden_files": false,
    "date_format": "relative"
  },
  "cleaner": {
    "safe_period_days": 1,
    "exclude_paths": [],
    "auto_clean_enabled": false
  },
  "security": {
    "secure_folder_enabled": false,
    "biometric_required": true
  }
}
\end{lstlisting}

\appendix

\chapter{Complexity Reference Card}

\begin{longtable}{@{}p{52mm}p{28mm}p{34mm}p{26mm}@{}}
\toprule
\textbf{Operation} & \textbf{Complexity} & \textbf{Latency (1M files)} & \textbf{RAM Tier} \\
\midrule
Exact path lookup & \bigO{1} & $<$25\,ns & HOT \\
Filename prefix search & \bigO{k} & $<$0.5\,ms & HOT \\
Single keyword search & \bigO{1} & $<$1\,ms & HOT \\
Multi-keyword AND & \bigO{N/64} & $<$2\,ms & HOT \\
Folder listing & \bigO{1} & $<$1\,ms & HOT \\
Filter by MIME type & \bigO{1} & $<$1\,ms & HOT \\
Range by file size & \bigO{\log\log U} & $<$1\,ms & WARM \\
Range by date & \bigO{\log\log U} & $<$1\,ms & WARM \\
Duplicate check & \bigO{1} & $<$1\,ms & WARM \\
Semantic AI search & \bigO{\log n} & $<$10\,ms & COLD (mmap) \\
Content search (PDF/DOCX) & \bigO{1} & $<$1\,ms & HOT \\
Logical delete & \bigO{N/64} & $<$2\,ms per 1k & HOT \\
Insert new file & \bigO{k + t} & $<$1\,ms & HOT \\
Rename / move & \bigO{1} & $<$0.5\,ms & HOT \\
Undo delete & \bigO{N/64} & $<$2\,ms per 1k & HOT \\
Recent files list & \bigO{1} & $<$0.1\,ms & HOT \\
Cross-app sync (Layer 1) & \bigO{1} & $<$1\,ms & HOT \\
Storage analytics (cached) & \bigO{1} & $<$0.1\,ms & HOT \\
Junk scan (full) & \bigO{n} & $<$3\,s & HOT \\
\bottomrule
\end{longtable}

\chapter{Data Structure Glossary}

\begin{description}[leftmargin=0pt, style=nextline]
  \item[\textbf{Radix Trie (Compressed Prefix Tree)}]
    A tree where each node stores a compressed string segment. Enables prefix matching in \bigO{k}$ time where $k$ is the query length, independent of collection size.

  \item[\textbf{RoaringBitmap}]
    A compressed bitset using three internal container types selected per data density. Achieves both space efficiency and fast set operations (\bigO{N/64} AND/OR).

  \item[\textbf{Van Emde Boas Tree (VEB)}]
    A tree for integer keys supporting predecessor, successor, and range queries in \bigO{\log\log U} time, where $U$ is the universe size.

  \item[\textbf{HNSW (Hierarchical Navigable Small World)}]
    A multi-layer graph structure for approximate nearest-neighbour search in vector spaces. Achieves \bigO{\log n} search with 95--98\% recall accuracy.

  \item[\textbf{xxHash64}]
    A non-cryptographic hash function with $\sim$10\,GB/s throughput. Used for path hashing (PathMap keys) and content checksums (ChecksumMap keys).

  \item[\textbf{Write-Ahead Log (WAL)}]
    A crash-safety mechanism where every index mutation is recorded in a sequential log before the in-memory index is modified. On crash recovery, the WAL is replayed.

  \item[\textbf{inotify}]
    A Linux kernel subsystem providing event notifications for filesystem changes. FileObserver on Android wraps inotify. Delivers change events in $<$1\,ms.

  \item[\textbf{MiniLM-L6}]
    A compact sentence embedding model producing 384-dimensional vectors. 22\,MB on disk, $\sim$15\,ms inference on modern Android SoC.

  \item[\textbf{Memory-Mapped File (mmap)}]
    A file mapped into the process virtual address space. The OS loads only accessed pages on demand --- enabling 600\,MB+ data structures without loading them fully into RAM.

  \item[\textbf{Tombstone / Logical Delete}]
    A deletion strategy where a deleted item is marked in a deletion BitSet rather than removed from all data structures immediately. User sees the item gone at \bigO{1}; physical cleanup is asynchronous.
\end{description}

\chapter{Android Permission Reference}

\begin{longtable}{@{}p{64mm}p{22mm}p{55mm}@{}}
\toprule
\textbf{Permission} & \textbf{API Level} & \textbf{FLUX Rationale} \\
\midrule
\texttt{READ\_EXTERNAL\_STORAGE} & 16--32 & Legacy storage read access \\
\texttt{WRITE\_EXTERNAL\_STORAGE} & 16--28 & Legacy storage write access \\
\texttt{MANAGE\_EXTERNAL\_STORAGE} & 30+ & Full filesystem access for complete index build \\
\texttt{READ\_MEDIA\_IMAGES} & 33+ & Read images (API 33+ replacement) \\
\texttt{READ\_MEDIA\_VIDEO} & 33+ & Read video files \\
\texttt{READ\_MEDIA\_AUDIO} & 33+ & Read audio files \\
\texttt{FOREGROUND\_SERVICE} & 9+ & Background index service notification \\
\texttt{RECEIVE\_BOOT\_COMPLETED} & all & Resume watching after device restart \\
\texttt{REQUEST\_IGNORE\_BATTERY\_OPTIMIZATIONS} & 23+ & Prevent WorkManager deferral \\
\texttt{USE\_BIOMETRIC} & 28+ & Secure folder authentication \\
\bottomrule
\end{longtable}
